% course 4, page 1

\bigskip

\begin{center}
\fbox{\textbf{ALGEBRA I}}\\[2pt]
(Curs 4)
\end{center}
\begin{flushright}
\textit{20 Oct.\ 1994}
\end{flushright}

\bigskip
\noindent\textbf{(1)} Lege de compoziție indusă pe o parte stabilă. $(A,\varphi)$, $\varphi : A\times A \to A$, $A\neq\varnothing$.

\noindent $B\subseteq A$, $B$ parte stab.\ a lui $A$ în raport cu $\varphi$ dacă $(\forall)\,x,y\in B \Rightarrow \varphi(x,y)\in B$.

\noindent $\varphi' : B\times B \to B$.

\bigskip
\noindent\textbf{(2)} Lege de compoziție indusă pe o mulțime factor. $(A,R)$

\noindent $\varphi : A\times A \to A$, \quad $R\subseteq A\times A$

\noindent $R$ compatibil cu $\varphi$: \quad $xRx'$ şi $yRy'$ $\Rightarrow$ $xyRx'y'$ (în notație multiplicativă.)

\[
\frac{A}{R} = \left\{\hat a \mid a\in A\right\} \qquad \varphi^* : \frac{A}{R}\times\frac{A}{R} \to \frac{A}{R} \qquad \varphi^*(\hat a,\hat b) \overset{\text{def}}{=} \widehat{\varphi(a,b)}
\]
\[
\hat a\hat b = \widehat{ab}; \qquad \hat a + \hat b = \widehat{a+b}.
\]

\bigskip
\noindent\textbf{(3)} Lege de compoziție indusă pe o mult.\ de funcții. $(A,\varphi)$, $I\neq\varnothing$.

\noindent $A^I = \{f \mid f:I\to A\}$

\noindent $\tilde\varphi$ = legea de comp.\ indusă de $\varphi$.

\noindent $(A^I,\tilde\varphi)$ \qquad $(\mathbb{R},+)$; $I=[a,b]$.

\noindent $\mathbb{R}^{[a,b]} = \{f \mid f:[a,b]\to\mathbb{R}\}$; \qquad $\mathbb{R}^{\{1,\ldots,n\}} = \{f \mid f:\{1,\ldots,n\}\to\mathbb{R}\}$.

\medskip
\noindent $A = a_0,a_1,\ldots,a_n,\ldots$

\noindent $m,n\in\mathbb{N}^*$.

\[
\mathbb{R}^{\{1,\ldots,m\}\times\{1,\ldots,n\}}
\]
\noindent $A : \{1,\ldots,m\}\times\{1,\ldots,n\} \to \mathbb{R}$ \hfill $\mathbb{R}\ni a_{ij} = A(i,j)$

\medskip
\noindent $(A,\varphi)$, $I\neq\varnothing$

\noindent $f,g\in A^I$

\[
f*g : I\to A \qquad (f*g)(x) \overset{\text{def}}{=} f(x)*g(x) \in A.
\]
\[
\varphi^* : A^I\times A^I \to A^I \qquad \big(\varphi^*(f,g)\big)(x) \overset{\text{def}}{=} \varphi\big(f(x),g(x)\big).
\]

\medskip
\noindent (1) Fie $f,g : [a,b]\to\mathbb{R}$; \quad $f+g : [a,b]\to\mathbb{R}$.

\noindent $(\forall)\, x\in[a,b]$ \qquad $(f+g)(x) = f(x)+g(x)$.

\noindent $fg : [a,b]\to\mathbb{R}$ \qquad $(fg)(x) = f(x)g(x)$.






% course 4, page 2
\section*{Curs 4 - Pagina 2}

\noindent\fbox{Teoremă} \ Fie $(A,\varphi)$; $\varphi : A\times A \to A$, $I\neq\varnothing$.

\[
\tilde\varphi : A^I \times A^I \to A^I
\]

\begin{enumerate}
\item[(1)] Dacă $\varphi$ asoc.\ (comut.) $\Rightarrow$ $\tilde\varphi$ asoc.\ (comut.).
\item[(2)] Dacă $e\in A$ el.\ neutru pt.\ $\varphi$ $\Rightarrow$ $\tilde e : I\to A$, $\tilde e(i)=e$ este el.\ neutru pt.\ $\tilde\varphi$.
\item[(3)] $\varphi$ are el.\ neutru şi $f(i)$ $(\forall)\,i$ e simetrizabil față de $\varphi$ $\Rightarrow$ $f$ simetrizabil față de $\tilde\varphi$ şi un simetric al lui $f$ este $f'(i) = (f(i))'$.
\end{enumerate}

\bigskip
\noindent (1) Notăm legea pe $A$ cu $*$.\footnote{formularea exactă e neclară în original (compactată/greu de citit) -- se pare că se redenumește legea $\varphi$ prin $*$ pentru comoditatea notației în demonstrație.}

\[
(f*g)(i) \overset{\text{def}}{=} f(i)*g(i) = g(i)*f(i) \overset{\text{def}}{=} (g*f)(i), \quad \text{adevărat.}
\]

\noindent (2) $(f*\tilde e)(i) = f(i)*\tilde e(i) = f(i)$.

\noindent $f*\tilde e = \tilde e*f = f$.

\noindent (3) $(f*f')(i) = \tilde e(i) = e$.

\noindent $f(i)*f'(i) = f(i)*(f(i))' = e$.

\bigskip
\noindent\textbf{(4)} Produs a 2 legi de compoziție.

\noindent $(A_1,\varphi_1)$, $(A_2,\varphi_2)$; $A = A_1\times A_2$.

\noindent $\varphi : A\times A \to A$; \quad $\varphi : (A_1\times A_2)\times(A_1\times A_2) \to A_1\times A_2$

\[
\varphi\big((x_1,x_2),(y_1,y_2)\big) \overset{\text{def}}{=} \big(\varphi_1(x_1,y_1),\varphi_2(x_2,y_2)\big).
\]
\[
(x_1,x_2)+(y_1,y_2) = (x_1+y_1,\,x_2+y_2) \quad \text{-- notație aditivă.}
\]

\noindent $(\mathbb{Z},+)$ \qquad $(\mathbb{Z}_6,+)$ \qquad $\mathbb{Z}\times\mathbb{Z}_6$

\bigskip
\noindent\fbox{Teoremă} \ Fie $(A_1,\varphi_1)$, $(A_2,\varphi_2)$ şi $\varphi = \varphi_1\times\varphi_2$.

\begin{enumerate}
\item[(1)] Dacă $\varphi_1$ şi $\varphi_2$ sunt asoc.\ (comut.) $\Rightarrow$ $\varphi$ asoc.\ (comut.)
\item[(2)] $e_1,e_2$ = elem.\ neutru pt.\ $\varphi_1,\varphi_2$ $\Rightarrow$ $(e_1,e_2)$ = el.\ neutru pt.\ $\varphi$.
\item[(3)] $x_1\in A_1$ e simetrizabil în rap.\ cu $\varphi_1$ şi $x_2\in A_2$ e sim.\ în rap.\ cu $\varphi_2$ $\Rightarrow$ $(x_1,x_2)$ e simetrizabil în rap.\ cu $\varphi$ şi $(x_1,x_2)' = (x_1',x_2')$.
\end{enumerate}





% course 4, page 3
\section*{Curs 4 - Pagina 3}
\noindent\textbf{(5)} Transportul unei legi de compoziție printr-o bijecție

\[
(A,\varphi) \qquad A \xrightarrow{f} A'
\]

\noindent Dacă $x',y'\in A' \Rightarrow \exists!\, x,y\in A$ a.î.\ $x'=f(x)$ şi $y'=f(y)$

\[
\varphi' : A'\times A' \to A', \qquad \varphi'(x',y') = f\big(\varphi(x,y)\big) \in A'
\]

\noindent [notație: $\varphi(x,y)=xy$] \quad $x'y' \overset{\text{def}}{=} f(xy)$\footnote{această linie e foarte compactată în original; se pare că se introduce convenția multiplicativă $\varphi(x,y)=xy$, astfel încât definiția lui $\varphi'$ devine $x'y'=f(xy)$.}

\[
f : A\to A' \qquad f(xy) = f(x)f(y) \quad (\forall)\,x,y\in A.
\]
\[
f^{-1}(x'y') = f^{-1}(x')f^{-1}(y') \quad (\forall)\,x',y'\in A'.
\]

\bigskip
\noindent\underline{Teoremă:} Fie $(A,\varphi)$, $\varphi:A\times A\to A$.

\noindent $f:A\to A'$ şi $\varphi':A'\times A'\to A'$.

\begin{enumerate}
\item[1)] $\varphi$ = asoc.\ (comut.) $\Leftrightarrow$ $\varphi'$ e asoc.\ (comut.).
\item[2)] $e\in A$ elem.\ neutru pt.\ $\varphi$ $\Leftrightarrow$ $f(e)$ el.\ neutru pt.\ $\varphi'$.\footnote{în original apare ,,pt.\ $f$'' / ,,pt.\ $f$' '' -- corectat aici la $\varphi$/$\varphi'$ pentru consecvență cu punctul 1) şi cu restul teoremei.}
\item[3)] $x\in A$ simetrizabil în raport cu $\varphi$ $\Leftrightarrow$ $f(x)$ simetrizabil în raport cu $\varphi'$.
\end{enumerate}

\noindent $x',y'\in A'$.

\[
x'y' = f(x)f(y) = f(xy) = f(yx) = f(y)f(x).
\]

\bigskip
\begin{center}
\fbox{\Large\textbf{GRUPURI}}
\end{center}

\noindent\underline{\textbf{Subgrupuri. T.\ lui Lagrange}}

\bigskip
\noindent\underline{Def:} Fie $G$ un grup. O submulțime $H$ nevidă a lui $G$ s.n.\ \textbf{subgrup} al grupului $G$ dacă:

\begin{enumerate}
\item[(1)] $(\forall)\,x,y\in H \Rightarrow xy\in H$
\item[(2)] $(\forall)\,x\in H \Rightarrow x^{-1}\in H$
\end{enumerate}

\bigskip
\noindent\underline{Pb:} Dat un grup $G$ să determinăm toate congruențele pe $G$.

\noindent Rel.\ de ech.\ pe $G$ compatibilă cu operația grupului:

\[
xRy \Rightarrow zxRzy.
\]
